
\lesson{19/10}
\subsubsection{Recap of Covariants}
\begin{itemize}
    \item $\bar\psi\psi$ is a scalar;
    \item $\bar\psi\gamma_5\psi$, since $\gamma_5$ anticommutes with
        $\gamma^\mu$, it commutes with $S$ and it is a scalar, but it changes
        sign under parity (${\gamma^0,\gamma_5}=0$), so this is a
        \textit{pseudo}scalar;
    \item $\bar\psi\gamma^\mu\psi$ is a Lorentz vector: it is the Noether current
        associated to the phase invariance;
    \item $\bar\psi\gamma^\mu\gamma_5\psi$ is a Lorentz pseudovector: it is the
        Noether current associated to the chiral phase invariance;
    \item $\bar\psi\sigma^{\mu\nu}\psi$ is a twice-covariant antisymmetric
        Lorentz tensor.
\end{itemize}
The 5 quantities between $\bar\psi$ and $\psi$, which are
$$\mathbb{I},\;\gamma_5,\;\gamma^\mu,\;\gamma^\mu\gamma_5,\;\sigma^{\mu\nu}$$
are collectively denoted as $\Gamma^a$.

\subsubsection{More on Noether currents}
Consider the continuity equation for $J^\mu$:
$$\pd_\mu J^\mu = \pd_\mu\bar\psi\gamma^\mu\psi + \bar\psi\gamma^\mu\pd_\mu\psi
$$
substituting from the two Dirac equations yields that the continuity equation
is always identically satisfied.

Consider now the continuity equation for $J^\mu_5:$
$$\pd_\mu J^\mu_5 = \pd_\mu\bar\psi\gamma^\mu\gamma^5\psi + \bar\psi
    \gamma^\mu\gamma_5\pd_\mu\psi$$
now to substitute from the equations of motion a minus sign pops out of the
anticommutation of $\gamma^\mu$, $\gamma_5$ and the continuity equation holds
\textbf{if and only if} the particle is massless.

\subsection{Chirality and Chiral projectors}
The chiral projectors are the operators
$$P_L = \frac{1}{2}(\mathbb{I}-\gamma_5),\; P_R=\frac{1}{2}(\mathbb{I}+\gamma_5)$$
and their effect is
$$P_L\psi = -\begin{pmatrix}
\psi_L\\0
\end{pmatrix},\;P_R \psi= \begin{pmatrix}
0\\\psi_R
\end{pmatrix}.
$$
\textbf{Def.} The chirality of a spinor is its eigenvalue under $\gamma_5$.
\subsection{Massless Dirac equation}
Since the chiral phase transformation is a symmetry if and only if $m=0$,
the massless Dirac equation must be considered.

Consider first the Dirac equation in matrix form
$$
\begin{pmatrix}
    -m\mathbb{I}&i(\pd_0+\sigma^i\pd_i)\\
    i(\pd_0-\sigma^i\pd_i)&-m\mathbb{I}
\end{pmatrix}\begin{pmatrix}
\psi_L\\\psi_R
\end{pmatrix}
=0
$$
and this can be expressed as 2 coupled equations
$$\begin{dcases*}
i\left(\pd_0 + \sigma^i\pd_i\right)\psi_R = m\psi_L\\
i\left(\pd_0 - \sigma^i\pd_i\right)\psi_L = m\psi_R
\end{dcases*}:
$$
the coupling constant is the mass, so the Dirac equation is reduced
to two uncoupled equations for the left and right spinors \textit{if and only if}
the particle is massless.

This means that massive fermions oscillate between left and right states, while
massless fermions remain in their \textit{chirality}.

Recalling that, for massless particles, $E  =|\mathbf{p}| = \pd_0$
and $p_i = \pd_i$, the equations are
in the form
$$\frac{\mathbf{\sigma \cdot p}}{|\mathbf{p}|}\psi_R = \psi_R$$
$$\frac{\mathbf{\sigma\cdot p}}{|\mathbf{p}|}\psi_L =- \psi_L:$$
the left and right spinors are eigenstates of the helicity operator
with eigenvalues $\pm1$.
